\documentclass[11pt,a4paper]{article}

\usepackage[utf8]{inputenc}
\usepackage[T1]{fontenc}
\usepackage[ngerman]{babel}
\usepackage[margin=2.5cm]{geometry}
\usepackage{listings}
\usepackage{xcolor}
\usepackage{enumitem}
\usepackage{booktabs}
\usepackage{array}
\usepackage{hyperref}
\usepackage{parskip}
\usepackage{titlesec}

\definecolor{codebg}{RGB}{242,242,242}
\definecolor{codefg}{RGB}{40,40,40}

\lstset{
  backgroundcolor=\color{codebg},
  basicstyle=\ttfamily\small\color{codefg},
  breaklines=true,
  frame=single,
  framerule=0pt,
  xleftmargin=4pt,
  xrightmargin=4pt,
  framesep=6pt,
  columns=fullflexible,
  showstringspaces=false,
}

\titleformat{\section}{\Large\bfseries}{\thesection}{1em}{}
\titleformat{\subsection}{\large\bfseries}{\thesubsection}{1em}{}

\hypersetup{
  colorlinks=true,
  linkcolor=blue!50!black,
  urlcolor=blue!50!black,
}

\title{\textbf{ROBO\_NEU -- Ablauf f\"ur verteilte BO-Benchmarks\\auf der LRZ Compute Cloud}}
\author{}
\date{Stand: 05.08.2026}

\begin{document}
\maketitle
\vspace{-1.5em}
\noindent\hrulefill

\section{\"Uberblick}

\texttt{ROBO\_NEU} rechnet BO-Benchmarks (Bayessche Optimierung auf BBOB/COCO-Funktionen)
\"uber mehrere unabh\"angige LRZ-Compute-Cloud-VMs parallel. Jede VM verarbeitet ihren
eigenen Task-Chunk mit mehreren Worker-Prozessen. Es gibt \textbf{keine zentrale
Steuerung} -- jede VM wird einzeln angesprochen.

\begin{table}[h]
\centering
\begin{tabular}{p{3.2cm} p{5.2cm} p{5.2cm}}
\toprule
\textbf{Begriff} & \textbf{Bedeutung} & \textbf{Hinweis} \\
\midrule
VM = OpenStack Instance & 1 VM entspricht immer genau 1 Instance & Anzahl \"uber LRZ-Kontingent (Quota) begrenzt \\
\texttt{instance\_0 ... instance\_N} & Nur Ordnernamen f\"ur Task-Chunks in ROBO\_NEU & Kein Bezug zum OpenStack-Begriff \\
Worker & Parallele Prozesse pro VM & Richtwert: 1 Worker pro vCPU \\
\bottomrule
\end{tabular}
\end{table}

\section{Ablaufdiagramm}

\begin{lstlisting}
Dein PC (prepare_tasks.py, split_tasks.py)
      |
      +-- scp instance_0 --> VM 0 -- SSH -- run_instance.py --workers N
      +-- scp instance_1 --> VM 1 -- SSH -- run_instance.py --workers N
      +-- scp instance_2 --> VM 2 -- SSH -- run_instance.py --workers N
      +-- scp instance_3 --> VM 3 -- SSH -- run_instance.py --workers N
                                                                    |
Dein PC <-- scp status.json / results.csv zurueck von jeder VM  <--+
      |
   merge_results.py --> results_gesamt.xlsx
\end{lstlisting}

\section{Schritt-f\"ur-Schritt-Anleitung}

\subsection{Schritt 1 -- Einmalig zentral (lokal oder auf VM 0)}

Erzeugt die vollst\"andige Task-Liste und teilt sie in Chunks auf, einen pro VM.

\begin{lstlisting}[language=bash]
python prepare_tasks.py
python split_tasks.py --instanzen 4
\end{lstlisting}

Ergebnis: \texttt{data/instances/instance\_0} \dots \texttt{instance\_3}, jeweils mit
\texttt{tasks.csv} (eigener Anteil) und \texttt{tensor\_noise.npy} (identische Kopie).

\subsection{Schritt 2 -- Pro VM: Code + Daten hochladen}

F\"ur jede VM einmal wiederholen (Code z.\,B. per \texttt{git clone}, Daten per
\texttt{scp}/\texttt{rsync}):

\begin{lstlisting}[language=bash]
scp -r data/instances/instance_0 user@vm0-ip:~/robo_neu/
scp -r data/instances/instance_1 user@vm1-ip:~/robo_neu/
scp -r data/instances/instance_2 user@vm2-ip:~/robo_neu/
scp -r data/instances/instance_3 user@vm3-ip:~/robo_neu/
\end{lstlisting}

\subsection{Schritt 3 -- Pro VM: einloggen und Prozess starten}

Jede VM bekommt eine eigene SSH-Session. \texttt{tmux} sorgt daf\"ur, dass der Prozess
bei SSH-Abbruch weiterl\"auft:

\begin{lstlisting}[language=bash]
ssh user@vm0-ip
cd ~/robo_neu
tmux new -s robo
python run_instance.py --instance-dir instance_0 --workers 10
# Strg+B, D zum Detachen der tmux-Session
\end{lstlisting}

Dasselbe wiederholen f\"ur \texttt{vm1/instance\_1}, \texttt{vm2/instance\_2},
\texttt{vm3/instance\_3} -- jede VM arbeitet unabh\"angig, ohne Wissen von den anderen.

\subsection{Schritt 4 -- Fortschritt pr\"ufen (jederzeit w\"ahrend des Laufs m\"oglich)}

\begin{lstlisting}[language=bash]
scp user@vm0-ip:~/robo_neu/instance_0/results/status.json data/instances/instance_0/results/
# ... fuer alle VMs wiederholen
python check_progress.py --pattern "data/instances/instance_*/results/status.json"
\end{lstlisting}

\subsection{Schritt 5 -- Nach Abschluss: Ergebnisse zusammenf\"uhren}

\begin{lstlisting}[language=bash]
scp user@vm0-ip:~/robo_neu/instance_0/results/results.csv data/instances/instance_0/results/
# ... fuer alle VMs wiederholen
python merge_results.py --pattern "data/instances/instance_*/results/results.csv" --out data/results_gesamt.xlsx
\end{lstlisting}

\section{Wichtige Hinweise}

\begin{itemize}[leftmargin=1.2em]
  \item Bei Absturz oder SSH-Abbruch: denselben \texttt{run\_instance.py}-Befehl einfach
        erneut ausf\"uhren -- bereits erledigte Tasks (\texttt{status=ok} in
        \texttt{results.csv}) werden automatisch \"ubersprungen.
  \item Fehlerhafte einzelne Tasks stoppen den Gesamtlauf nicht
        (\texttt{status="error"} wird protokolliert, Rest l\"auft weiter).
  \item Worker-Anzahl sollte der Anzahl vCPUs der jeweiligen VM entsprechen -- mehr
        Worker als Kerne bringt keinen Geschwindigkeitsvorteil.
  \item Mehr VMs als aktuell 4 sind m\"oglich, erfordern aber ggf. eine
        Quota-Erh\"ohung beim LRZ (vCPU-/Instanz-Kontingent des
        Compute-Cloud-Projekts).
  \item \texttt{config.py} ist aktuell auf Testwerte reduziert -- f\"ur den echten
        Lauf (500.000+ Kombinationen) \texttt{FUNKTIONEN}, \texttt{DIMENSIONEN},
        \texttt{SURROGATE\_AKQUISITION\_PAARE} und \texttt{N\_SAMPLE\_INIT} wieder
        hochskalieren.
\end{itemize}

\section{Offene Automatisierungsm\"oglichkeit (optional)}

Da Schritt 2--4 pro VM identisch sind, l\"asst sich der Ablauf mit einem kleinen
Wrapper-Skript (Bash/PowerShell) \"uber eine Liste von VM-IPs automatisieren, statt
jede VM manuell einzeln anzusteuern -- auf Wunsch kann ein solches Skript erg\"anzt
werden.

\end{document}
